\ifmostrarGabarito
\begin{minted}{python}
import pandas as pd
from typing import List, Dict, Any

def filtrar_prontidao_navios(
    dados_frota: List[Dict[str, Any]],
    autonomia_min: float,
    missoes_alvo: List[str]
) -> pd.DataFrame:
    """
    Cria DataFrame de frota, aplica filtros com .loc[] e ordena por prontidão de autonomia.
    """
    df = pd.DataFrame(dados_frota)
    
    # Máscara booleana composta com operador bitwise & e método .isin()
    filtro = (df["autonomia_dias"] >= autonomia_min) & (df["missao_designada"].isin(missoes_alvo))
    
    # Seleção explícita de linhas e colunas com .loc
    colunas_interesse = ["nome", "classe", "autonomia_dias"]
    df_selecionado = df.loc[filtro, colunas_interesse].copy()
    
    # Ordenação decrescente e redefinição de índice
    df_final = df_selecionado.sort_values(
        by="autonomia_dias", ascending=False
    ).reset_index(drop=True)
    
    return df_final
\end{minted}

\textbf{Explicação:}
A função constrói o DataFrame tabular a partir de registros brutos e aplica filtragem indexada com \texttt{.loc[linhas, colunas]}. Os filtros lógicos combinam expressões com parênteses obrigatórios e o operador \texttt{\&}. A ordenação com \texttt{.sort\_values()} e \texttt{.reset\_index(drop=True)} produz uma saída limpa e padronizada.
\fi
